% GENERATED FILE. Edit canonical lessons, then run npm run generate:books.
% canonical-source-hash: fnv1a64:aeaefcce6b1457a9
% canonical-lessons: GE-C459-bauernhof, GE-C459-weide, GE-C459-scheune, GE-C459-grasen, GE-C459-melken

\chapter{Auf dem Land}
\label{ch:ge-auf-dem-land}
\hlchaptermodality{459}

\begin{chapteropening}
\textbf{By the end of this chapter:} \emph{I can talk about a farm: the farm, the pasture, the barn, grazing and milking.}

Complete the last lesson of chapter 459: I can talk about a farm: the farm, the pasture, the barn, grazing and milking.
\end{chapteropening}

\section[der Bauernhof]{der Bauernhof — the farm}

\label{lesson:GE-C459-bauernhof}



\begin{quote}
Before the new one: say the German for the stag, then the German for to chirp.
\end{quote}

\subsection*{You'll want to know: der Bauernhof}
\textbf{der Bauernhof} — \textquotedblleft{}the farm\textquotedblright{}.

\begin{quote}
Im Urlaub wohnen wir auf einem Bauernhof.
\end{quote}

\begin{quote}
\emph{On holiday we stay on a farm.}
\end{quote}

Plural: \emph{die Bauernhöfe}.

\begin{cousinweb}
\emph{Bauernhof} is built from \emph{Bauer} 'farmer', inherited from Old High German \emph{būr} (EtymDB), and \emph{Hof} 'yard, farm'.
\end{cousinweb}

\begin{grammarlens}[title={What you've built}]
One more word for telling, explaining and describing at B1.
\end{grammarlens}

\subsection*{Your turn}
\begin{itemize}
\raggedright
  \item \emph{Say it:} \emph{der Bauernhof}
  \item \emph{Say it:} \emph{der Bauernhof}, once more
  \item \emph{Say it:} a sentence of your own with \emph{der Bauernhof}
  \item \emph{Recall:} say the German for the stag, then the German for to chirp, then say \emph{der Bauernhof} again
\end{itemize}

\begin{tcolorbox}[breakable,colback=teal!4,colframe=teal!35!black,title={Before you move on}]
What does der Bauernhof mean? (\textquotedblleft{}The farm\textquotedblright{}.) Say it once more.
\end{tcolorbox}
\section[die Weide]{die Weide — the pasture (for animals); the willow}

\label{lesson:GE-C459-weide}



\begin{quote}
Before the new one: say the German for to chirp, then the German for the farm.
\end{quote}

\subsection*{You'll want to know: die Weide}
\textbf{die Weide} — \textquotedblleft{}the pasture (for animals); the willow\textquotedblright{}.

\begin{quote}
Die Kühe stehen den ganzen Tag auf der Weide.
\end{quote}

\begin{quote}
\emph{The cows stand in the pasture all day.}
\end{quote}

Plural: \emph{die Weiden}. \emph{Die Weide} is also a willow tree, a different word.

\begin{cousinweb}
\emph{Weide} 'pasture' is inherited from Old High German \emph{weida} (EtymDB). \emph{Weide} 'willow' is a separate inherited word, Old High German \emph{wīda} (EtymDB).
\end{cousinweb}

\begin{grammarlens}[title={What you've built}]
One more word for telling, explaining and describing at B1.
\end{grammarlens}

\subsection*{Your turn}
\begin{itemize}
\raggedright
  \item \emph{Say it:} \emph{die Weide}
  \item \emph{Say it:} \emph{die Weide}, once more
  \item \emph{Say it:} a sentence of your own with \emph{die Weide}
  \item \emph{Recall:} say the German for to chirp, then the German for the farm, then say \emph{die Weide} again
\end{itemize}

\begin{tcolorbox}[breakable,colback=teal!4,colframe=teal!35!black,title={Before you move on}]
What does die Weide mean? (\textquotedblleft{}The pasture (for animals); the willow\textquotedblright{}.) Say it once more.
\end{tcolorbox}
\section[die Scheune]{die Scheune — the barn}

\label{lesson:GE-C459-scheune}



\begin{quote}
Before the new one: say the German for the farm, then the German for the pasture.
\end{quote}

\subsection*{You'll want to know: die Scheune}
\textbf{die Scheune} — \textquotedblleft{}the barn\textquotedblright{}.

\begin{quote}
In der Scheune liegt das Heu für den Winter.
\end{quote}

\begin{quote}
\emph{The hay for the winter is in the barn.}
\end{quote}

Plural: \emph{die Scheunen}.

\begin{cousinweb}
\emph{Scheune} is inherited through Middle High German \emph{schiune} from Old High German \emph{scugina} 'shed' (EtymDB).
\end{cousinweb}

\begin{grammarlens}[title={What you've built}]
One more word for telling, explaining and describing at B1.
\end{grammarlens}

\subsection*{Your turn}
\begin{itemize}
\raggedright
  \item \emph{Say it:} \emph{die Scheune}
  \item \emph{Say it:} \emph{die Scheune}, once more
  \item \emph{Say it:} a sentence of your own with \emph{die Scheune}
  \item \emph{Recall:} say the German for the farm, then the German for the pasture, then say \emph{die Scheune} again
\end{itemize}

\begin{tcolorbox}[breakable,colback=teal!4,colframe=teal!35!black,title={Before you move on}]
What does die Scheune mean? (\textquotedblleft{}The barn\textquotedblright{}.) Say it once more.
\end{tcolorbox}
\section[grasen]{grasen — to graze}

\label{lesson:GE-C459-grasen}



\begin{quote}
Before the new one: say the German for the pasture, then the German for the barn.
\end{quote}

\subsection*{You'll want to know: grasen}
\textbf{grasen} — \textquotedblleft{}to graze\textquotedblright{}.

\begin{quote}
Die Schafe grasen auf der Wiese.
\end{quote}

\begin{quote}
\emph{The sheep are grazing in the meadow.}
\end{quote}

\emph{Es grast}, \emph{sie haben gegrast}.

\begin{cousinweb}
\emph{grasen} is built on \emph{Gras} 'grass', inherited from Proto-Germanic \emph{grasą} (EtymDB lists the verb).
\end{cousinweb}

\begin{grammarlens}[title={What you've built}]
One more word for telling, explaining and describing at B1.
\end{grammarlens}

\subsection*{Your turn}
\begin{itemize}
\raggedright
  \item \emph{Say it:} \emph{grasen}
  \item \emph{Say it:} \emph{grasen}, once more
  \item \emph{Say it:} a sentence of your own with \emph{grasen}
  \item \emph{Recall:} say the German for the pasture, then the German for the barn, then say \emph{grasen} again
\end{itemize}

\begin{tcolorbox}[breakable,colback=teal!4,colframe=teal!35!black,title={Before you move on}]
What does grasen mean? (\textquotedblleft{}To graze\textquotedblright{}.) Say it once more.
\end{tcolorbox}
\section[melken]{melken — to milk}

\label{lesson:GE-C459-melken}



\begin{quote}
Before the new one: say the German for the barn, then the German for to graze.
\end{quote}

\subsection*{You'll want to know: melken}
\textbf{melken} — \textquotedblleft{}to milk\textquotedblright{}.

\begin{quote}
Der Bauer melkt die Kühe um fünf Uhr morgens.
\end{quote}

\begin{quote}
\emph{The farmer milks the cows at five in the morning.}
\end{quote}

Past: \emph{melkte}, \emph{hat gemolken}.

\begin{cousinweb}
\emph{melken} is inherited from Old High German \emph{melchan} and Proto-Germanic \emph{melkaną} (EtymDB); English \emph{milk} is a cousin.
\end{cousinweb}

\begin{grammarlens}[title={What you've built}]
One more word for telling, explaining and describing at B1.
\end{grammarlens}

\subsection*{Your turn}
\begin{itemize}
\raggedright
  \item \emph{Say it:} \emph{melken}
  \item \emph{Say it:} \emph{melken}, once more
  \item \emph{Say it:} a sentence of your own with \emph{melken}
  \item \emph{Recall:} say the German for the barn, then the German for to graze, then say \emph{melken} again
\end{itemize}

\begin{tcolorbox}[breakable,colback=teal!4,colframe=teal!35!black,title={Before you move on}]
What does melken mean? (\textquotedblleft{}To milk\textquotedblright{}.) Say it once more.
\end{tcolorbox}
